\documentclass[10pt,a4paper]{amsart}
\usepackage{amsmath,amssymb,amsthm,graphicx}
% Retained minimal preamble for the Stacks Project chapter flat.tex (More on Flatness).
% Verbatim excerpt of the chapter preamble: Xy-pic setup (native preamble lines 14-18) and the native
% theorem declarations with the mathematical macros (native preamble lines 154-217). Copied mechanically;
% nothing added, removed or redefined. The wrapper class and the amsmath/amssymb/amsthm/graphicx base
% packages are supplied by the assembler, so the class/externaldocument part of the native preamble is omitted.
% For commutative diagrams we use Xy-pic
\usepackage[all]{xy}

% We use 2cell for 2-commutative diagrams.
\xyoption{2cell}
% Theorem environments.
%
\theoremstyle{plain}
\newtheorem{theorem}[subsection]{Theorem}
\newtheorem{proposition}[subsection]{Proposition}
\newtheorem{lemma}[subsection]{Lemma}

\theoremstyle{definition}
\newtheorem{definition}[subsection]{Definition}
\newtheorem{example}[subsection]{Example}
\newtheorem{exercise}[subsection]{Exercise}
\newtheorem{situation}[subsection]{Situation}

\theoremstyle{remark}
\newtheorem{remark}[subsection]{Remark}
\newtheorem{remarks}[subsection]{Remarks}

\numberwithin{equation}{subsection}

% Macros
%
\def\lim{\mathop{\mathrm{lim}}\nolimits}
\def\colim{\mathop{\mathrm{colim}}\nolimits}
\def\Spec{\mathop{\mathrm{Spec}}}
\def\Hom{\mathop{\mathrm{Hom}}\nolimits}
\def\Ext{\mathop{\mathrm{Ext}}\nolimits}
\def\SheafHom{\mathop{\mathcal{H}\!\mathit{om}}\nolimits}
\def\SheafExt{\mathop{\mathcal{E}\!\mathit{xt}}\nolimits}
\def\Sch{\mathit{Sch}}
\def\Mor{\mathop{\mathrm{Mor}}\nolimits}
\def\Ob{\mathop{\mathrm{Ob}}\nolimits}
\def\Sh{\mathop{\mathit{Sh}}\nolimits}
\def\NL{\mathop{N\!L}\nolimits}
\def\CH{\mathop{\mathrm{CH}}\nolimits}
\def\proetale{{pro\text{-}\acute{e}tale}}
\def\etale{{\acute{e}tale}}
\def\QCoh{\mathit{QCoh}}
\def\Ker{\mathop{\mathrm{Ker}}}
\def\Im{\mathop{\mathrm{Im}}}
\def\Coker{\mathop{\mathrm{Coker}}}
\def\Coim{\mathop{\mathrm{Coim}}}

% Boxtimes
%
\DeclareMathSymbol{\boxtimes}{\mathbin}{AMSa}{"02}

%
% Macros for moduli stacks/spaces
%
\def\QCohstack{\mathcal{QC}\!\mathit{oh}}
\def\Cohstack{\mathcal{C}\!\mathit{oh}}
\def\Spacesstack{\mathcal{S}\!\mathit{paces}}
\def\Quotfunctor{\mathrm{Quot}}
\def\Hilbfunctor{\mathrm{Hilb}}
\def\Curvesstack{\mathcal{C}\!\mathit{urves}}
\def\Polarizedstack{\mathcal{P}\!\mathit{olarized}}
\def\Complexesstack{\mathcal{C}\!\mathit{omplexes}}
% \Pic is the operator that assigns to X its picard group, usage \Pic(X)
% \Picardstack_{X/B} denotes the Picard stack of X over B
% \Picardfunctor_{X/B} denotes the Picard functor of X over B
\def\Pic{\mathop{\mathrm{Pic}}\nolimits}
\def\Picardstack{\mathcal{P}\!\mathit{ic}}
\def\Picardfunctor{\mathrm{Pic}}
\def\Deformationcategory{\mathcal{D}\!\mathit{ef}}

% --- appended 2026-10-12 (Chief): the project own definitions of the reference/slogan/history
% display environments, copied mechanically from the shared Stacks preamble lines 8-12
% (\usepackage{verbatim} plus the three \newenvironment lines, which the chapter preamble keeps
% verbatim from the project preamble). Nothing added, removed or redefined: these environments
% carry no mathematical content and the project itself typesets them as comments. Without these
% three lines a counted statement whose own text uses them fails with Environment slogan undefined.
\usepackage{verbatim}
\newenvironment{reference}{\comment}{\endcomment}
\newenvironment{slogan}{\comment}{\endcomment}
\newenvironment{history}{\comment}{\endcomment}
\usepackage{cleveref}
\crefname{originallemma}{original lemma}{original lemmas}
\Crefname{originallemma}{Original lemma}{Original lemmas}
\makeatletter
\newlabel{algebra-lemma-map-between-finitely-presented}{{05GF}{source}{}{sourcepointer-1}{}}
\newlabel{algebra-lemma-map-between-finitely-presented@cref}{{[originallemma][0][]05GF}{[1][0][]source}{}{}{}}
\makeatother
\crefname{originallemma}{original lemma}{original lemmas}
\Crefname{originallemma}{Original lemma}{Original lemmas}
\makeatletter
\newlabel{algebra-lemma-universally-injective-check-stalks}{{05CL}{source}{}{sourcepointer-2}{}}
\newlabel{algebra-lemma-universally-injective-check-stalks@cref}{{[originallemma][0][]05CL}{[1][0][]source}{}{}{}}
\makeatother
\crefname{originallemma}{original lemma}{original lemmas}
\Crefname{originallemma}{Original lemma}{Original lemmas}
\makeatletter
\newlabel{algebra-proposition-fppf-open}{{00I1}{source}{}{sourcepointer-3}{}}
\newlabel{algebra-proposition-fppf-open@cref}{{[originallemma][0][]00I1}{[1][0][]source}{}{}{}}
\makeatother
\crefname{originallemma}{original lemma}{original lemmas}
\Crefname{originallemma}{Original lemma}{Original lemmas}
\makeatletter
\newlabel{algebra-theorem-openness-flatness}{{00RC}{source}{}{sourcepointer-4}{}}
\newlabel{algebra-theorem-openness-flatness@cref}{{[originallemma][0][]00RC}{[1][0][]source}{}{}{}}
\makeatother
\crefname{originallemma}{original lemma}{original lemmas}
\Crefname{originallemma}{Original lemma}{Original lemmas}
\makeatletter
\newlabel{lemma-induction-step}{{05I3}{source}{}{sourcepointer-5}{}}
\newlabel{lemma-induction-step@cref}{{[originallemma][0][]05I3}{[1][0][]source}{}{}{}}
\makeatother
\title{The induction step for a finitely presented flat ring map with geometrically integral fibres}
\author{The Stacks Project Authors}
\date{stacks-project:flat.tex:local-snapshot-20261011; GNU FDL 1.2 or later}
\begin{document}
\maketitle
\noindent\textit{Editorial note.} All delivered mathematical text is the author's own, copied from the pinned Stacks Project chapter file flat.tex; no statement or proof step was rewritten, reordered or removed.\par\medskip
\noindent\textit{Editorial note.} Preamble provenance: the chapter's own retained minimal preamble (sources/stacks-flat/retained-preamble.tex, lines 1-85) is a verbatim excerpt of the native Xy-pic setup and of the native theorem declarations and mathematical macros, followed by the project's own reference/slogan/history display-environment definitions copied mechanically from the shared preamble; the wrapper class and the amsmath/amssymb/amsthm/graphicx base packages are supplied by the assembler, and extra\_packages is left empty because everything needed is already in the retained preamble. No macro was invented or redefined.\par\medskip
\noindent\textit{Editorial note.} Same-chapter pointers are located at their \textbackslash{}label lines inside the pinned chapter file; cross-chapter pointers are located in the project's own tag list staged in the snapshot (sources/stacks-chow/tags.txt), at the verbatim lines carrying their canonical tags.\par\medskip
\noindent\textit{Allowed standard prerequisites (editorial).}\begin{itemize}
\item The algebraic induction step for a map of modules over such a ring map. Same chapter flat.tex, label lemma-induction-step, invoked at proof line 3181. Applied to the localised map alpha at r, it gives the equivalence of conditions (1), (3), (4), (5) and (6); the remaining implication (1) implies (2) is what the counted proof supplies.
\item Openness of flatness, openness of the locus where a map of finitely presented modules is an isomorphism, and openness of a flat finitely presented ring map. Algebra, labels algebra-theorem-openness-flatness (proof line 3190), algebra-lemma-map-between-finitely-presented (proof line 3203) and algebra-proposition-fppf-open (proof line 3207). Makes the flat locus U\_1, the isomorphism locus U\_2 and the image Phi(U\_1 intersect U\_2) open in the appropriate spectra, which is what allows the choice of f in R and then of g in S.
\item Testing universal injectivity on stalks. Algebra, label algebra-lemma-universally-injective-check-stalks, invoked at proof line 3218. Converts universal injectivity of alpha at every prime of the open set Phi(U\_1 intersect U\_2) into universal injectivity of the localised map alpha\_f, which is the condition required by item (2).
\end{itemize}\medskip
\ifdefined\hypertarget\hypertarget{sourcepointer-1}{}\fi
\noindent\textit{Source reference.} Original lemma 05GF: sources/stacks-chow/tags.txt, source lines 7059--7059. Navigation pointer for \textbackslash{}ref\{algebra-lemma-map-between-finitely-presented\} used inside the counted span (line 3203); the argument that consumes it is retained in the span.\par\medskip
\ifdefined\hypertarget\hypertarget{sourcepointer-2}{}\fi
\noindent\textit{Source reference.} Original lemma 05CL: sources/stacks-chow/tags.txt, source lines 6924--6924. Navigation pointer for \textbackslash{}ref\{algebra-lemma-universally-injective-check-stalks\} used inside the counted span (line 3218); the argument that consumes it is retained in the span.\par\medskip
\ifdefined\hypertarget\hypertarget{sourcepointer-3}{}\fi
\noindent\textit{Source reference.} Original lemma 00I1: sources/stacks-chow/tags.txt, source lines 663--663. Navigation pointer for \textbackslash{}ref\{algebra-proposition-fppf-open\} used inside the counted span (line 3207); the argument that consumes it is retained in the span.\par\medskip
\ifdefined\hypertarget\hypertarget{sourcepointer-4}{}\fi
\noindent\textit{Source reference.} Original lemma 00RC: sources/stacks-chow/tags.txt, source lines 998--998. Navigation pointer for \textbackslash{}ref\{algebra-theorem-openness-flatness\} used inside the counted span (line 3190); the argument that consumes it is retained in the span.\par\medskip
\ifdefined\hypertarget\hypertarget{sourcepointer-5}{}\fi
\noindent\textit{Source reference.} Original lemma 05I3: sources/stacks-flat/pinned-original-source.tex, source lines 2356--2356. Navigation pointer for \textbackslash{}ref\{lemma-induction-step\} used inside the counted span (line 3181); the argument that consumes it is retained in the span.\par\medskip
\begin{lemma}
\label{lemma-induction-step-fp}
Let $R$ be a ring. Let $R \to S$ be a finitely presented
flat ring map with geometrically integral fibres. Let
$\mathfrak q \subset S$ be a prime ideal lying over the prime
$\mathfrak r \subset R$. Set $\mathfrak p = \mathfrak r S$.
Let $N$ be a finitely presented $S$-module.
There exists $r \geq 0$ and an $S$-module map
$$
\alpha : S^{\oplus r} \longrightarrow N
$$
such that
$\alpha : \kappa(\mathfrak p)^{\oplus r} \to N \otimes_S \kappa(\mathfrak p)$
is an isomorphism. For any such $\alpha$ the following are equivalent:
\begin{enumerate}
\item $N_{\mathfrak q}$ is $R$-flat,
\item there exists an $f \in R$, $f \not \in \mathfrak r$ such that
$\alpha_f : S_f^{\oplus r} \to N_f$ is $R_f$-universally injective and
a $g \in S$, $g \not \in \mathfrak q$ such that $\Coker(\alpha)_g$
is $R$-flat,
\item $\alpha_{\mathfrak r}$ is $R_{\mathfrak r}$-universally injective and
$\Coker(\alpha)_{\mathfrak q}$ is $R$-flat
\item $\alpha_{\mathfrak r}$ is injective and
$\Coker(\alpha)_{\mathfrak q}$ is $R$-flat,
\item $\alpha_{\mathfrak p}$ is an isomorphism and
$\Coker(\alpha)_{\mathfrak q}$ is $R$-flat, and
\item $\alpha_{\mathfrak q}$ is injective and
$\Coker(\alpha)_{\mathfrak q}$ is $R$-flat.
\end{enumerate}
\end{lemma}

\begin{proof}
To obtain $\alpha$ set
$r = \dim_{\kappa(\mathfrak p)} N \otimes_S \kappa(\mathfrak p)$ and pick
$x_1, \ldots, x_r \in N$ which form a basis of
$N \otimes_S \kappa(\mathfrak p)$. Define
$\alpha(s_1, \ldots, s_r) = \sum s_i x_i$. This proves the existence.

\medskip\noindent
Fix a choice of $\alpha$.
We may apply
Lemma \ref{lemma-induction-step}
to the map
$\alpha_{\mathfrak r} : S_{\mathfrak r}^{\oplus r} \to N_{\mathfrak r}$.
Hence we see that (1), (3), (4), (5), and (6) are all equivalent.
Since it is also clear that (2) implies (3) we see that all we have to
do is show that (1) implies (2).

\medskip\noindent
Assume (1). By openness of flatness, see
Algebra, Theorem \ref{algebra-theorem-openness-flatness},
the set
$$
U_1 = \{\mathfrak q' \subset S \mid N_{\mathfrak q'}\text{ is flat over }R\}
$$
is open in $\Spec(S)$. It contains $\mathfrak q$ by assumption
and hence $\mathfrak p$. Because $S^{\oplus r}$ and $N$ are finitely presented
$S$-modules the set
$$
U_2 = \{\mathfrak q' \subset S \mid
\alpha_{\mathfrak q'}\text{ is an isomorphism}\}
$$
is open in $\Spec(S)$, see
Algebra, Lemma \ref{algebra-lemma-map-between-finitely-presented}.
It contains $\mathfrak p$ by (5). As $R \to S$
is finitely presented and flat the map
$\Phi : \Spec(S) \to \Spec(R)$ is open, see
Algebra, Proposition \ref{algebra-proposition-fppf-open}.
For any prime $\mathfrak r' \in \Phi(U_1 \cap U_2)$ we see that
there exists a prime $\mathfrak q'$ lying over $\mathfrak r'$ such that
$N_{\mathfrak q'}$ is flat and such that $\alpha_{\mathfrak q'}$ is
an isomorphism, which implies that $\alpha \otimes \kappa(\mathfrak p')$
is an isomorphism where $\mathfrak p' = \mathfrak r' S$. Thus
$\alpha_{\mathfrak r'}$ is $R_{\mathfrak r'}$-universally injective
by the implication (1) $\Rightarrow$ (3).
Hence if we pick $f \in R$, $f \not \in \mathfrak r$ such that
$D(f) \subset \Phi(U_1 \cap U_2)$ then we conclude that
$\alpha_f$ is $R_f$-universally injective, see
Algebra, Lemma \ref{algebra-lemma-universally-injective-check-stalks}.
The same reasoning also shows that for any
$\mathfrak q' \in U_1 \cap \Phi^{-1}(\Phi(U_1 \cap U_2))$
the module $\Coker(\alpha)_{\mathfrak q'}$ is $R$-flat.
Note that $\mathfrak q \in U_1 \cap \Phi^{-1}(\Phi(U_1 \cap U_2))$.
Hence we can find a $g \in S$, $g \not \in \mathfrak q$ such
that $D(g) \subset U_1 \cap \Phi^{-1}(\Phi(U_1 \cap U_2))$
and we win.
\end{proof}

\end{document}
